\chapter{Chapter06 — K Wise Independent Lower Bound}
%%% generated by the blueprint pipeline from FABL.Chapter06.Constructions.KWiseIndependentLowerBound
\section{Overview}

Declarations of the Lean module \texttt{FABL.Chapter06.Constructions.KWiseIndependentLowerBound}.

\section{Declarations}

\begin{definition}[Density of a sign multiset]
\lean{FABL.signMultisetDensity}
\statementsource{Analysis-of-Boolean-Functions-by-Ryan-ODonnell}{pdf-ffbe9c7b3306-p144-b007}
\label{FABL.signMultisetDensity}
\leanok
\uses{FABL.ProbabilityDensity, FABL.ProbabilityDensity.uniformPushforward, FABL.binaryCubeSignEquiv}
Let $\Omega$ be a finite nonempty set and let $A$ assign to each element of $\Omega$ a
sign vector in $\{-1,1\}^n$. The associated probability density on $\mathbb{F}_2^n$ is
the density of the uniform pushforward of $\Omega$ under the map that converts each sign
vector to its corresponding binary vector.
\end{definition}

\begin{definition}[Multiset k-wise independence]
\lean{FABL.IsKWiseIndependentMultiset}
\statementsource{Analysis-of-Boolean-Functions-by-Ryan-ODonnell}{
  pdf-ffbe9c7b3306-p146-b004,
  pdf-ffbe9c7b3306-p147-b006,
  pdf-ffbe9c7b3306-p170-b003
}
\label{FABL.IsKWiseIndependentMultiset}
\leanok
\uses{FABL.IsLowDegreeFourierRegular, FABL.binaryFunctionOnSignCube, FABL.signMultisetDensity}
Let $\Omega$ be a finite nonempty set, let $n,k$ be nonnegative integers, and let $A$
assign to each element of $\Omega$ a sign vector in ${\pm1}^n$. The multiset given by
these vectors is called $k$-wise independent when the binary function associated with
its sign-multiset density is low-degree Fourier-regular with parameters $0$ and $k$.
\end{definition}

\begin{theorem}[Fourier coefficient of an empirical sign distribution]
\lean{FABL.fourierCoeff_signMultisetDensity}
\label{FABL.fourierCoeff_signMultisetDensity}
\leanok
\uses{FABL.ProbabilityDensity.expectation, FABL.ProbabilityDensity.expectation_uniformPushforward, FABL.binaryFunctionOnSignCube, FABL.fourierCoeff, FABL.fourierCoeff_binaryFunctionOnSignCube, FABL.monomial, FABL.monomial_binaryCubeSignEquiv, FABL.signMultisetDensity, FABL.χ}
\difficulty{3}
Let $n$ be a nonnegative integer, let $\Omega$ be a finite nonempty set, let $A$ assign
to each element of $\Omega$ a point of the $n$-dimensional sign cube, and let
$S\subseteq[n]$. The Fourier coefficient at $S$ of the sign-cube density induced by the
uniform distribution on $\Omega$ equals the uniform average over $\Omega$ of the parity
character on $S$ evaluated at $A(\omega)$.
\end{theorem}

\begin{theorem}[Orthogonality of monomials under limited independence]
\lean{FABL.IsKWiseIndependentMultiset.expect_monomial_mul}
\label{FABL.IsKWiseIndependentMultiset.expect_monomial_mul}
\leanok
\uses{FABL.IsKWiseIndependentMultiset, FABL.binaryFunctionOnSignCube, FABL.fourierCoeff, FABL.fourierCoeff_signMultisetDensity, FABL.monomial, FABL.monomial_mul_monomial, FABL.signMultisetDensity}
\difficulty{5}
Let $n,k\ge0$, let $\Omega$ be a finite nonempty set, and let $A$ assign to each
$\omega\in\Omega$ a sign vector in $\{-1,1\}^n$. For a subset
$S\subseteq\{1,\ldots,n\}$, write $\chi_S(x)=\prod_{i\in S}x_i$. Suppose the
distribution of $A(\omega)$, with $\omega$ uniform on $\Omega$, has all nonconstant
Fourier coefficients indexed by sets of size at most $k$ equal to zero. Then for any
subsets $S,T\subseteq\{1,\ldots,n\}$ with $|S\cup T|\le k$, the uniform expectation over
$\omega\in\Omega$ of $\chi_S(A(\omega))\chi_T(A(\omega))$ is $1$ if $S=T$ and $0$
otherwise.
\end{theorem}

\begin{definition}[Normalized restricted Walsh vector]
\lean{FABL.normalizedRestrictedWalshVector}
\label{FABL.normalizedRestrictedWalshVector}
\leanok
\uses{FABL.monomial}
Let $\Omega$ be a finite nonempty set, let $A$ assign to each $\omega\in\Omega$ a sign
vector in $\{ -1,1\}^n$, and let $S\subseteq\{0,\ldots,n-1\}$. The normalized restricted
Walsh vector is the function on $\Omega$ given by $\omega\mapsto
|\Omega|^{-1/2}\prod_{i\in S}A(\omega)_i$, regarded as a vector in Euclidean space.
\end{definition}

\begin{theorem}[Orthonormality of restricted Walsh vectors]
\lean{FABL.normalizedRestrictedWalshVector_orthonormal}
\label{FABL.normalizedRestrictedWalshVector_orthonormal}
\leanok
\uses{FABL.IsKWiseIndependentMultiset, FABL.IsKWiseIndependentMultiset.expect_monomial_mul, FABL.monomial, FABL.normalizedRestrictedWalshVector}
\difficulty{5}
Let $\Omega$ be a finite nonempty set, let $A$ assign a sign vector in $\{−1,1\}^n$ to
each element of $\Omega$, and suppose the multiset of vectors given by $A$ is $k$-wise
independent. Let $\mathcal F$ be a finite family of subsets of $\{1,\ldots,n\}$ such
that $|S\cup T|\le k$ for every $S,T\in\mathcal F$. For each $S\in\mathcal F$, form the
vector indexed by $\omega\in\Omega$ whose entry is $|\Omega|^{-1/2}\prod_{i\in S}
A(\omega)_i$; these vectors are orthonormal in the Euclidean space of real-valued
functions on $\Omega$.
\end{theorem}

\begin{theorem}[Bound on a frequency family]
\lean{FABL.card_frequencyFamily_le_card_of_isKWiseIndependentMultiset}
\label{FABL.card_frequencyFamily_le_card_of_isKWiseIndependentMultiset}
\leanok
\uses{FABL.IsKWiseIndependentMultiset, FABL.normalizedRestrictedWalshVector_orthonormal}
\difficulty{2}
Let $\Omega$ be a finite nonempty set, and assign to each $\omega\in\Omega$ a sign
vector on $n$ coordinates. Suppose the resulting multiset of sign vectors is $k$-wise
independent, in the sense that every nonempty Fourier coefficient of degree at most $k$
of its induced density is zero. If $\mathcal F$ is a family of coordinate subsets such
that $|S\cup T|\leq k$ for every $S,T\in\mathcal F$, then $|\mathcal F|\leq|\Omega|$.
\end{theorem}

\begin{theorem}[Union size of low-degree sets]
\lean{FABL.card_union_le_two_mul_of_mem_lowDegreeFourierFamily}
\label{FABL.card_union_le_two_mul_of_mem_lowDegreeFourierFamily}
\leanok
\uses{FABL.lowDegreeFourierFamily, FABL.mem_lowDegreeFourierFamily}
\difficulty{3}
Let $n,m$ be natural numbers, and let $S,T$ be subsets of an $n$-element set. If each of
$S$ and $T$ has at most $m$ elements, then $|S\cup T|\le 2m$.
\end{theorem}

\begin{theorem}[Even-wise frequency family]
\lean{FABL.exists_even_kWiseFrequencyFamily}
\label{FABL.exists_even_kWiseFrequencyFamily}
\leanok
\uses{FABL.card_lowDegreeFourierFamily_eq_sum_choose, FABL.card_union_le_two_mul_of_mem_lowDegreeFourierFamily, FABL.lowDegreeFourierFamily}
\difficulty{3}
Let $n$ and $k$ be nonnegative integers, with $k$ even. There is a family $\mathcal F$
of subsets of an $n$-element set such that $|S\cup T|\le k$ for every $S,T\in\mathcal
F$, and $|\mathcal F|=\sum_{j=0}^{k/2}\binom{n}{j}$.
\end{theorem}

\begin{definition}[Odd-size frequency extension]
\lean{FABL.oddKWiseFrequencyExtension}
\label{FABL.oddKWiseFrequencyExtension}
\leanok
For natural numbers $n,m$ and an index $i\in\{0,\ldots,n-1\}$, form the family of sets
by taking every $m$-element subset of the indices other than $i$ and adjoining $i$ to
it.
\end{definition}

\begin{theorem}[Cardinality of an odd frequency extension member]
\lean{FABL.mem_oddKWiseFrequencyExtension_card}
\label{FABL.mem_oddKWiseFrequencyExtension_card}
\leanok
\uses{FABL.oddKWiseFrequencyExtension}
\difficulty{4}
Let $n,m$ be natural numbers, let $i$ be an index in an $n$-element set, and let $S$
belong to the family formed by adjoining $i$ to each $m$-element subset of the other
indices. Then $i\in S$ and $|S|=m+1$.
\end{theorem}

\begin{theorem}[Cardinality of the odd frequency extension]
\lean{FABL.card_oddKWiseFrequencyExtension}
\label{FABL.card_oddKWiseFrequencyExtension}
\leanok
\uses{FABL.oddKWiseFrequencyExtension}
\difficulty{5}
For natural numbers $n$ and $m$ and a coordinate $i\in\{0,\ldots,n-1\}$, form the family
of sets obtained by adjoining $i$ to each $m$-element subset of the other $n-1$
coordinates. This family has cardinality $\binom{n-1}{m}$.
\end{theorem}

\begin{definition}[Odd-size frequency family]
\lean{FABL.oddKWiseFrequencyFamily}
\label{FABL.oddKWiseFrequencyFamily}
\leanok
\uses{FABL.lowDegreeFourierFamily, FABL.oddKWiseFrequencyExtension}
For natural numbers $n,m$ and an index $i\in\{0,\ldots,n-1\}$, the odd-size frequency
family is the union of the low-degree Fourier family with parameters $n,m$ and the
family obtained by adjoining $i$ to each $m$-element subset of the other indices.
\end{definition}

\begin{theorem}[Disjoint low-degree and extended frequencies]
\lean{FABL.disjoint_lowDegreeFourierFamily_oddKWiseFrequencyExtension}
\label{FABL.disjoint_lowDegreeFourierFamily_oddKWiseFrequencyExtension}
\leanok
\uses{FABL.lowDegreeFourierFamily, FABL.mem_lowDegreeFourierFamily, FABL.mem_oddKWiseFrequencyExtension_card, FABL.oddKWiseFrequencyExtension}
\difficulty{3}
Let $n,m$ be nonnegative integers and let $i$ be a coordinate in $\{0,\ldots,n-1\}$. The
family of all subsets of the coordinate set of size at most $m$ is disjoint from the
family formed by adjoining $i$ to each $m$-element subset of the other coordinates.
\end{theorem}

\begin{theorem}[Cardinality of the odd frequency family]
\lean{FABL.card_oddKWiseFrequencyFamily}
\label{FABL.card_oddKWiseFrequencyFamily}
\leanok
\uses{FABL.card_lowDegreeFourierFamily_eq_sum_choose, FABL.card_oddKWiseFrequencyExtension, FABL.disjoint_lowDegreeFourierFamily_oddKWiseFrequencyExtension, FABL.oddKWiseFrequencyFamily}
\difficulty{1}
For natural numbers $n$ and $m$ and a coordinate $i\in\{0,\ldots,n-1\}$, the family
formed by taking all subsets of the $n$ coordinates of size at most $m$, together with
all sets obtained by adjoining $i$ to an $m$-element subset of the other coordinates,
has cardinality $\displaystyle \sum_{j=0}^{m}\binom{n}{j}+\binom{n-1}{m}$.
\end{theorem}

\begin{theorem}[Union size in a frequency family]
\lean{FABL.card_union_le_two_mul_add_one_of_mem_oddKWiseFrequencyFamily}
\label{FABL.card_union_le_two_mul_add_one_of_mem_oddKWiseFrequencyFamily}
\leanok
\uses{FABL.card_union_le_two_mul_of_mem_lowDegreeFourierFamily, FABL.mem_lowDegreeFourierFamily, FABL.mem_oddKWiseFrequencyExtension_card, FABL.oddKWiseFrequencyFamily}
\difficulty{5}
Let $n$ and $m$ be natural numbers, let $i$ be a coordinate among $n$ coordinates, and
let $S$ and $T$ belong to the following family: all subsets of size at most $m$,
together with all subsets of size $m+1$ that contain $i$. Then $|S\cup T|\le 2m+1$.
\end{theorem}

\begin{theorem}[Odd-wise frequency family size]
\lean{FABL.exists_odd_kWiseFrequencyFamily}
\label{FABL.exists_odd_kWiseFrequencyFamily}
\leanok
\uses{FABL.card_oddKWiseFrequencyFamily, FABL.card_union_le_two_mul_add_one_of_mem_oddKWiseFrequencyFamily, FABL.oddKWiseFrequencyFamily}
\difficulty{3}
Let $n$ be a positive integer and let $k$ be odd. There is a family $\\mathcal F$ of
subsets of $[n]$ such that $|S\\cup T|\\le k$ for every $S,T\\in\\mathcal F$, and
\[|\\mathcal F|=\\sum_{j=0}^{(k-1)/2}\\binom{n}{j}+\\binom{n-1}{(k-1)/2}.\]
\end{theorem}

\begin{theorem}[Even k wise independent multiset card lower bound]
\lean{FABL.even_kWiseIndependentMultiset_card_lowerBound}
\statementsource{Analysis-of-Boolean-Functions-by-Ryan-ODonnell}{pdf-ffbe9c7b3306-p170-b003}
\label{FABL.even_kWiseIndependentMultiset_card_lowerBound}
\leanok
\uses{FABL.IsKWiseIndependentMultiset, FABL.card_frequencyFamily_le_card_of_isKWiseIndependentMultiset, FABL.exists_even_kWiseFrequencyFamily}
\difficulty{3}
% dataset-pipeline: stub -- prose pending build_informal_statements.py
Statement pending regeneration.
\end{theorem}

\begin{theorem}[Cardinality bound for odd-wise independence]
\lean{FABL.odd_kWiseIndependentMultiset_card_lowerBound}
\proofsource{Analysis-of-Boolean-Functions-by-Ryan-ODonnell}{pdf-ffbe9c7b3306-p170-b003}
\label{FABL.odd_kWiseIndependentMultiset_card_lowerBound}
\leanok
\uses{FABL.IsKWiseIndependentMultiset, FABL.card_frequencyFamily_le_card_of_isKWiseIndependentMultiset, FABL.exists_odd_kWiseFrequencyFamily}
\difficulty{3}
Let $n>0$, let $k$ be an odd natural number, and let $A$ be a nonempty finite multiset
of sign vectors in $\{-1,1\}^n$ whose associated density has every nonconstant Fourier
coefficient of degree at most $k$ equal to zero. Then $\displaystyle
\sum_{j=0}^{(k-1)/2}\binom{n}{j}+\binom{n-1}{(k-1)/2}\le |A|$, where $|A|$ counts
multiset elements.
\end{theorem}

\begin{theorem}[Size bound for a wise independent multiset]
\lean{FABL.choose_floor_half_le_card_of_isKWiseIndependentMultiset}
\label{FABL.choose_floor_half_le_card_of_isKWiseIndependentMultiset}
\leanok
\uses{FABL.IsKWiseIndependentMultiset, FABL.card_frequencyFamily_le_card_of_isKWiseIndependentMultiset, FABL.card_union_le_two_mul_of_mem_lowDegreeFourierFamily, FABL.lowDegreeFourierFamily, FABL.mem_lowDegreeFourierFamily}
\difficulty{4}
Let $n$ and $k$ be natural numbers, let $\Omega$ be a finite nonempty set, and let $A$
assign a sign vector in $\{-1,1\}^{n}$ to each element of $\Omega$. If the resulting
multiset is $k$-wise independent, then $\binom{n}{\lfloor k/2\rfloor}\leq |\Omega|$.
\end{theorem}

\begin{theorem}[Size bound for independent multisets]
\lean{FABL.kWiseIndependentMultiset_card_isOmega}
\proofsource{Analysis-of-Boolean-Functions-by-Ryan-ODonnell}{pdf-ffbe9c7b3306-p170-b003}
\label{FABL.kWiseIndependentMultiset_card_isOmega}
\leanok
\uses{FABL.IsKWiseIndependentMultiset, FABL.choose_floor_half_le_card_of_isKWiseIndependentMultiset}
\difficulty{4}
Fix a nonnegative integer $k$. For each nonnegative integer $n$, let $\Omega_n$ be a
finite nonempty set and let $A_n$ map $\Omega_n$ into the $n$-dimensional sign cube.
Suppose that, for every $n$, the multiset of points given by $A_n$ has all nonempty
Fourier coefficients indexed by sets of at most $k$ coordinates equal to zero. Then, as
$n$ tends to infinity, $|\Omega_n|=O\!\left(n^{\lfloor k/2\rfloor}\right)$.
\end{theorem}
